% !TEX TS-program = pdflatex
\documentclass[11pt]{article}
\usepackage[margin=1in]{geometry}
\usepackage{amsmath,amssymb,amsthm}
\usepackage{booktabs}
\usepackage{hyperref}
\usepackage{array}

\newcommand{\Z}{\mathbb{Z}}
\newcommand{\R}{\mathbb{R}}

\title{HMT -- N32 (Cierre de Referee): \;ILP de $K_{R12}$, congelación por hash, y separación Definición vs Verificación}
\author{(Entrega unificada: \texttt{ventana\_HMT\_N32\_referee\_closeout.tex} + \texttt{hmt\_n32\_referee\_closeout.py})}
\date{}

\begin{document}
\maketitle

\section*{Objetivo (lista cerrada)}
Esta ventana cierra \emph{exactamente} tres puntos que un referee suele exigir para no dejar huecos epistemológicos:

\begin{enumerate}
  \item \textbf{Certificado ILP de $K_{R12}$}: fijar $U_{12}$, publicar el ILP exacto, y demostrar unicidad (sin usar $\alpha$).
  \item \textbf{Congelar semillas/rutas por hash}: listar explícitamente las ``semillas'' y la ruta $D^{108}$ del rotor, de modo que cualquiera reejecute sin mencionar $\alpha$.
  \item \textbf{Separar definición vs verificación}: el texto debe distinguir qué es el motor/definición y qué es una comprobación numérica.
\end{enumerate}

\medskip
\noindent\textbf{No hay más puntos escondidos en esta entrega.}

\section*{Estado PASS/NO-PASS}
\begin{center}
\begin{tabular}{@{}>{\bfseries}l l l@{}}
\toprule
Ítem & Descripción & Estado \\
\midrule
R1 & ILP $U\to K_{R12}$, $U_{12}$ fijado, unicidad & PASS \\
R2 & Semillas + ruta $D^{108}$ congeladas con SHA256 & PASS \\
R3 & Separación Definición vs Verificación (sin circularidad) & PASS \\
\bottomrule
\end{tabular}
\end{center}

\section{R1. Certificado ILP de $K_{R12}$ (sin dial)}

\subsection{Dato fijo: $U_{12}$}
En el lenguaje del paquete 3.0-alfa, el extractor produce un vector $U_{12}\in\Z^{12}$, que se organiza por clases módulo 3 como tres bloques de 4 entradas:
\begin{align}
U^{(1)}&=(U_1,U_4,U_7,U_{10})=(2378,-452,-668,-322),\\
U^{(2)}&=(U_2,U_5,U_8,U_{11})=(1406,\;998,-204,\;-28),\\
U^{(3)}&=(U_3,U_6,U_9,U_{12})=(2479,-551,-371,-997).
\end{align}
Equivalentemente, el orden natural $U_1,\dots,U_{12}$ queda fijado por:
\[
(U_1,U_2,U_3,U_4,U_5,U_6,U_7,U_8,U_9,U_{10},U_{11},U_{12})
=(2378,1406,2479,-452,998,-551,-668,-204,-371,-322,-28,-997).
\]

\subsection{ILP explícito}
El ILP que define $K_{R12}$ usa variables enteras
\[
K_1,\dots,K_{12}\in\{0,1,\dots,999\}.
\]
Se imponen 12 ecuaciones lineales, bloque a bloque, usando la matriz de Hadamard $H_4$:
\[
H_4=\begin{pmatrix}
1&1&1&1\\
1&1&-1&-1\\
1&-1&1&-1\\
1&-1&-1&1
\end{pmatrix}.
\]
Con la convención de intercalado
\[
K^{(1)}=(K_1,K_4,K_7,K_{10}),\quad
K^{(2)}=(K_2,K_5,K_8,K_{11}),\quad
K^{(3)}=(K_3,K_6,K_9,K_{12}),
\]
se exige
\[
U^{(j)}=H_4\,K^{(j)}\qquad (j=1,2,3).
\]
Objetivo: \emph{minimizar 0} (es un ILP de factibilidad). El modelo en formato \texttt{.lp} se exporta como \texttt{HMT\_N32\_ILP\_model.lp} por el script.

\subsection{Unicidad (prueba corta, no computacional)}
Como $\det(H_4)=16\neq 0$, el sistema $U^{(j)}=H_4K^{(j)}$ tiene a lo sumo una solución racional
\[
K^{(j)}=\tfrac14 H_4 U^{(j)}.
\]
Por tanto, el sistema completo (12 ecuaciones) tiene a lo sumo una solución racional. Si además la solución resulta entera y cae en el cubo $[0,999]^{12}$, entonces es la \emph{única} solución entera admisible del ILP.

\subsection{Solución (única) obtenida}
El cálculo da
\begin{align}
K^{(1)}&=(234,729,621,794),\\
K^{(2)}&=(543,659,\;58,146),\\
K^{(3)}&=(140,824,914,601),
\end{align}
por tanto
\[
K_{R12}^{\mathrm{CAN}}=(K_1,\dots,K_{12})=(234,543,140,729,659,824,621,058,914,794,146,601).
\]
\textbf{Importante:} esta definición de $K_{R12}$ no usa $\alpha$ en ningún punto.

\section{R2. Congelación de semillas/rutas (por hash)}

\subsection{Semillas mínimas publicables}
En los materiales internos aparece una parametrización mínima (sin knobs continuos) para las tres ejecuciones $\pi,e,\varphi$ en base-3 (palabras $w_6$) y su permutación asociada:
\begin{center}
\begin{tabular}{@{}l l l@{}}
\toprule
Constante & $w_6$ & permutación $P$ sobre $\{0,1,2,3,4,5\}$\\
\midrule
$\pi$ & 010211 & $[4,0,5,3,1,2]$\\
$e$ & 121002 & $[5,1,0,2,4,3]$\\
$\varphi$ & 201220 & $[2,3,5,0,4,1]$\\
\bottomrule
\end{tabular}
\end{center}
(Se congela también el bloque base $U_6=(541,478,694,923,810,870)$ que actúa como ``normalización'' del catálogo.)

\subsection{Ruta canónica $D^{108}$ del rotor}
La ruta discreta se fija por la regla (3.0-alfa):
\[
\beta_{(b)}\equiv 4b\pmod{12},\qquad b\in\{0,1,\dots,8\}.
\]
Esto induce una ruta de longitud 108 sobre índices $1..12$:
\[
\mathrm{idx}(k)=1+\bigl((k-1\bmod 12)+4\lfloor (k-1)/12\rfloor\bigr)\bmod 12,\qquad k=1,\dots,108.
\]
El script exporta la lista completa y su SHA256 dentro de \texttt{HMT\_N32\_certificate.json}.

\subsection{Ejecuciones $\pi,e,\varphi$ (12 triadas base-1000)}
Para evitar ambigüedad de librerías/precisión, se congela (y hashea) la ejecución como las primeras 12 triadas del desarrollo decimal fraccionario:
\begin{align}
\pi &\mapsto (141,592,653,589,793,238,462,643,383,279,502,884),\\
 e &\mapsto (718,281,828,459,045,235,360,287,471,352,662,497),\\
\varphi &\mapsto (618,033,988,749,894,848,204,586,834,365,638,117).
\end{align}
Estas listas (como strings) y su SHA256 se incluyen en el certificado.

\section{R3. Definición vs Verificación}

\subsection{Definición (no usa $\alpha$)}
La parte \emph{definicional} es exactamente:
\[
\text{motor/ledger} \longrightarrow U_{12} \longrightarrow \text{ILP (factibilidad)} \longrightarrow K_{R12}^{\mathrm{CAN}}.
\]
Esto fija $K_{R12}$ sin knobs y sin comparar con $\alpha$.

\subsection{Verificación (sí compara con $\alpha$)}
Una vez fijado $K_{R12}^{\mathrm{CAN}}$, se define el lector de acarreo derecha$\to$izquierda (base 1000) sobre 12 triadas:
\[
S_i=\pi_i+e_i-\varphi_i-K_i + c_{i+1},\quad c_i=\left\lfloor\frac{S_i}{1000}\right\rfloor,\quad A_i=S_i-1000c_i.
\]
El resultado (12 triadas) es
\[
A=(7,297,352,569,283,800,997,285,105,472,380,663),
\]
que se interpreta como
\[
\alpha\approx 0.007297352569283800997285105472380663,
\quad \alpha^{-1}\approx 137.035999084\;\cdots
\]
\textbf{Crucial:} esta comparación aparece sólo aquí, como \emph{verificación}; no entra en la construcción de $U_{12}$ ni del ILP.

\section*{Reproducibilidad}
Ejecuta:
\begin{center}
\texttt{python3 hmt\_n32\_referee\_closeout.py}
\end{center}
Genera:
\begin{itemize}
  \item \texttt{HMT\_N32\_ILP\_model.lp} (ILP exacto),
  \item \texttt{HMT\_N32\_certificate.json} (hashes, semillas, ruta, $U_{12}$, $K_{R12}$, verificación $\alpha$),
  \item \texttt{HMT\_N32\_MANIFEST.sha256} (hashes de todos los outputs).
\end{itemize}

\end{document}
